\section*{第 \ref{ch:b1:organism-environment} 章——生物体：与环境相互作用的系统}

\begin{solution}{exo:b1:organism-environment:1}
物质（哺乳动物摄食并呼出二氧化碳；植物吸收
二氧化碳并释放氧气），能量（哺乳动物摄取食物的化学能
并释放热量；植物吸收光并释放热量），
信息（哺乳动物看到捕食者；植物向光弯曲）。
\end{solution}

\begin{solution}{exo:b1:organism-environment:2}
$S/V = 3/r$ 减半。表面吸收量按 $r^2$ 增长，而
消耗量按 $r^3$ 增长；超过某一半径后，表面已无法满足
体积的需要，因此细胞只能保持在微米量级，或使其表面折叠。
\end{solution}

\begin{solution}{exo:b1:organism-environment:3}
通过\omterm{def:b1:organism-environment:homeostasis}{负反馈}，使\omterm{def:b1:organism-environment:milieu}{内环境}的某一变量维持在\omterm{def:b1:organism-environment:homeostasis}{调定点}附近。
感受器：皮肤和
脑中的温度感受器；整合中枢为下丘脑，其\omterm{def:b1:organism-environment:homeostasis}{调定点}为
\qty{37}{\celsius}；效应器：战栗的肌肉、皮肤血管、汗
腺。
\end{solution}

\begin{solution}{exo:b1:organism-environment:4}
蜥蜴：从 $1.0$ 降至 $0.25$，倍数为 $4$（两次加倍，
对应 \qty{20}{\celsius}）。小鼠：从 $1.0$ 升至 $3.4$，
倍数约为 $3.4$。
\end{solution}

\begin{solution}{exo:b1:organism-environment:5}
边长 \qty{1}{mm}：需要 $\qty{1}{nmol/s}$（$V = \qty{1}{mm^3}$）；供应
$2\times 6 = \qty{12}{nmol/s}$：足够。边长 \qty{10}{mm}：需要
\qty{1000}{nmol/s}，供应 $2\times 600 = \qty{1200}{nmol/s}$：勉强足够。
极限：由 $2\times 6L^2 = L^3$ 得 $L = \qty{12}{mm}$。
\end{solution}

\begin{solution}{exo:b1:organism-environment:6}
马：$3.4\times 500^{0.75} = \qty{360}{W}$，\qty{0.72}{W/kg}。尖鼠：
$3.4\times 0.005^{0.75} = \qty{0.064}{W}$，\qty{13}{W/kg}。每日：
$3\times 0.064\times 86400 = \qty{16.6}{kJ}$，即 \qty{3.3}{g} 食物，
约为其自身质量的三分之二。
\end{solution}

\begin{solution}{exo:b1:organism-environment:7}
$\qty{80}{W} = \qty{4.8}{kJ/min}$，即 \qty{0.24}{L} 的 $\mathrm{O_2}$
每分钟；每天 $80\times 86400 = \qty{6.9}{MJ}$。\qty{2000}{kcal}
= \qty{8.36}{MJ} 每天 = \qty{97}{W}。
\end{solution}

\begin{solution}{exo:b1:organism-environment:8}
小鼠：$100\times(0.02/70)^{0.75} = \qty{0.22}{m^2}$；大象：
$100\times(4000/70)^{0.75} = \qty{2080}{m^2}$。球体模型：小鼠
$r = \qty{1.7}{cm}$，$S = \qty{0.0036}{m^2}$（肺面积为
体表面积的六十倍）；大象 $r = \qty{0.98}{m}$，$S = \qty{12}{m^2}$（肺面积为
体表面积的一百七十倍）。动物越大，其肺
相对于外表面就越需要折叠。
\end{solution}

\begin{solution}{exo:b1:organism-environment:9}
效应器只有在感受器检测到偏差之后才起作用，因此
必须先出现偏差才能被纠正；纠正
随后会稍微过冲，感受器才报告其偏差。调节带宽度由
感受器的灵敏度（可检测的最小偏差）、
检测与效应之间的延迟，以及
效应器的强度决定。
\end{solution}

\begin{solution}{exo:b1:organism-environment:10}
$V = \qty{2e-5}{m^3}$，$r = \qty{1.68}{cm}$，$S = 4\pi r^2 =
\qty{3.6e-3}{m^2}$。散热损失：$5\times 3.6\times 10^{-3}\times 37 =
\qty{0.66}{W}$。基础产热为 \qty{0.18}{W}，约为其四分之一。
小鼠必须提高产热（战栗、进食：
达到基础值的三至四倍），降低 $h$（筑巢、聚集取暖）或让体温下降（蛰伏）。
\end{solution}

\begin{solution}{exo:b1:organism-environment:11}
散热 $\propto M^{2/3}$，产热 $\propto M^{3/4}$：比值
产热/散热 $\propto M^{1/12}$ 随质量增大，因此小动物
单位表面积的产热最少，大动物最多。球体
质量为 $M$：$S = 4\pi(3M/4000\pi)^{2/3} = \num{0.0483}\,M^{2/3}$（国际单位制）；
散热 $= 5\times 20\times 0.0483\,M^{2/3} = 4.83\,M^{2/3}$；令其与
$3.4\,M^{3/4}$ 相等，得 $M^{1/12} = 1.42$，$M \approx \qty{67}{kg}$。低于此质量，
静息的动物在无保温的情况下会感到寒冷；高于此质量，它会升温，
必须散热。
\end{solution}

\begin{solution}{exo:b1:organism-environment:12}
并非更简单，只是组织方式不同。植物的交换发生于外部
（叶和根的表面伸展到介质中），其细胞
暴露于其无法调节的介质中；其“调节”是生长与
形态，即调整表面本身。哺乳动物将其交换
转入受调节的体液之后，使其细胞从不接触
外界。\omterm{def:b1:organism-environment:homeostasis}{稳态}是动物\omterm{def:b1:organism-environment:milieu}{内环境}的性质，
而非复杂程度的度量：植物的细胞调节自身
内容物，其气孔调节水分流失，其身体是一个
不断扩大的交换表面——一种策略，而非缺失。
\end{solution}

\begin{solution}{pb:b1:organism-environment:1}
\textbf{1.} $V = M/\rho$：小鼠 \qty{2e-5}{m^3}（\qty{20}{cm^3}），
大象 \qty{4}{m^3}。
\textbf{2.} $r = (3V/4\pi)^{1/3}$：小鼠 \qty{1.68}{cm}，大象
\qty{0.985}{m}。
\textbf{3.} $S = 4\pi r^2$：小鼠 \qty{3.55e-3}{m^2}（\qty{35.5}{cm^2}），
大象 \qty{12.2}{m^2}。
\textbf{4.} $S/V = 3/r$：小鼠 \qty{178}{m^{-1}}，大象
\qty{3.05}{m^{-1}}；比值为 $58$。
\textbf{5.} $S/V \propto 1/r \propto M^{-1/3}$，因此比值为
$(M_e/M_m)^{1/3} = (2\times 10^5)^{1/3} = 58.5$。
\textbf{6.} 两者均更大（球体在给定体积下
表面积最小）；小鼠因尾、耳和细长的四肢而偏离
更多。
\textbf{7.} $P = 10\times S\times 20$：小鼠 \qty{0.71}{W}，大象
\qty{2440}{W}。
\textbf{8.} 小鼠 $3.4\times 0.02^{0.75} = \qty{0.18}{W}$；大象
$3.4\times 4000^{0.75} = \qty{1710}{W}$。
\textbf{9.} 散热 \qty{2440}{W}，静息时产热 \qty{1710}{W}，
裸露皮肤大体足够；大象在活动或处于较温暖空气中时，
有过热而非降温的危险。
\textbf{10.} 散热 \qty{0.71}{W}，相对 \qty{0.18}{W}，
多出约四倍。毛皮必须将 $h$ 降至约 \qty{2.5}{W.m^{-2}.K^{-1}}。
\textbf{11.} 其 $S/V$ 是所有陆生哺乳动物中最小的；耳朵
使其表面增加约三分之一，且因薄而富含血液，是散热器，
当身体产热超过 \qty{12}{m^2} 所能
散失的量时，将热量排散出去。
\textbf{12.} $S/V \propto M^{-1/3}$ 随 $M$ 减小而无限增大，
而 $B/M \propto M^{-1/4}$ 增长较慢；质量低于几克时，任何毛皮
或任何摄食速率都无法弥补散热。夜间让体温下降
（蛰伏）可去掉 $T_b - T_a$ 项，即在无食的数小时内。
\textbf{13.} $2B\times\qty{86400}{s}$：小鼠 \qty{31}{kJ}，大象
\qty{295}{MJ}。
\textbf{14.} 小鼠 $31/18 = \qty{1.7}{g}$；大象
$295\,000/5 = \qty{59}{kg}$。
\textbf{15.} 小鼠为 $1.7/20 = 8.6\%$（占其质量）；大象 $59/4000 =
1.5\%$：按每克计算，小鼠的进食量是大象的六倍。
\textbf{16.} 肠长之比 $35/0.4 = 88$；半径之比 $0.985/0.0168
= 59$。肠道的增长快于身体的线性尺度（大致
此处约为 $M^{0.36}$），这增加了大象的吸收表面
相对于其体积的比例。
\textbf{17.} $f_0 = 600\times 0.02^{1/4} = \qty{226}{min^{-1}}$；
$\tau_0 = 2/0.02^{1/4} = \qty{5.3}{yr}$。
\textbf{18.} $f = 226\times 4000^{-1/4} = \qty{28}{min^{-1}}$；$\tau =
5.3\times 4000^{1/4} = \qty{42}{yr}$。
\textbf{19.} 小鼠 $600\times 60\times 24\times 365\times 2 =
\num{6.3e8}$；大象 $28\times 60\times 24\times 365\times 42 =
\num{6.2e8}$。
\textbf{20.} 一生的心跳次数 $= f\tau \propto M^{-1/4}M^{1/4} = M^0$：
与质量无关，接近十亿次。
\textbf{21.} $f = 226\times 70^{-1/4} = \qty{78}{min^{-1}}$（接近
70）；$\tau = 5.3\times 70^{1/4} = \qty{15}{yr}$，远低于 80：
人类寿命是离群值，为哺乳动物趋势的三至五倍，
其原因（大脑、社会性、医学）不属于标度律的范畴。
\textbf{22.} $160\times 4000^{-1/4}/0.02^{-1/4} = 160\times 0.0473 =
\qty{7.6}{min^{-1}}$；每次呼吸的心跳数 $600/160 = 28/7.6 = 3.75$
两者相同，因为两个速率具有相同的指数。
\textbf{23.} $B/M = 3.4\,M^{3/4}/M = 3.4\,M^{-1/4}$；指数为 $-1/4$。
\textbf{24.} 小鼠 $0.18/0.02 = \qty{9.0}{W/kg}$；大象 $1710/4000 =
\qty{0.43}{W/kg}$。
\textbf{25.} 比值 $9.0/0.43 = 21$，等于 $(M_e/M_m)^{1/4} =
(2\times 10^5)^{1/4} = 21.1$：小鼠每克的代谢速率
是大象每克的二十一倍左右——单位质量\omterm{def:b1:organism-environment:allometry}{代谢率}之比
即为质量比的四次方根。
\end{solution}
